\section{Areal--volumetric incidence and topological entropy}

The TPK calendar induces the hinge between the additive and
multiplicative modes. In the primitive block
\[
 \tau_3=(+1,-1,0),
\]
the rule
\[
 +1:m\mapsto\Sigma,
 \qquad
 -1:m\mapsto\Pi,
 \qquad
 0:m\mapsto m
\]
produces the cyclic orbit \((\Sigma,\Pi,\Pi)\). The geometric reader
\[
 \Sigma\longmapsto\mathscr H_{\rm ar},
 \qquad
 \Pi\longmapsto\mathscr H_{\rm vol}
\]
transports these modes to the areal and volumetric sheets.

\begin{theorem}[Necessary incidence of the sheets]
\label{thm:incidencia-hojas-anexo}
In the order \((\mathscr H_{\rm ar},\mathscr H_{\rm vol})\), the
primitive incidence matrix is
\begin{equation}
 \boxed{
 F_{\rm av}=
 \begin{pmatrix}0&1\\1&1\end{pmatrix}.}
 \label{eq:fav-anexo-informacion}
\end{equation}
Calendars of lengths \(9,27,108\) yield
\[
 N_9=3F_{\rm av},
 \qquad N_{27}=9F_{\rm av},
 \qquad N_{108}=36F_{\rm av}.
\]
\end{theorem}

\begin{proof}
The consecutive pairs of the cycle are
\(\Sigma\to\Pi\), \(\Pi\to\Pi\), and \(\Pi\to\Sigma\). Each occurs
once, and \(\Sigma\to\Sigma\) does not occur. This fixes all four entries of
\(F_{\rm av}\). The longer calendars contain, respectively, three,
nine, and thirty-six copies of the primitive block.
\end{proof}

\begin{theorem}[Fibonacci law and measure of maximal entropy]
\label{thm:fibonacci-parry-anexo}
For \(n\geq1\),
\begin{equation}
 F_{\rm av}^{n}
 =\begin{pmatrix}F_{n-1}&F_n\\F_n&F_{n+1}\end{pmatrix}.
 \label{eq:potencias-fav-anexo}
\end{equation}
The one-step subshift defined by \(F_{\rm av}\) has
\begin{equation}
 \boxed{h_{\rm top}=\log\varphi},
 \qquad
 \boxed{\frac{\mu_{\rm ar}}{\mu_{\rm vol}}=\varphi^{-2}}.
 \label{eq:entropia-parry-hojas-anexo}
\end{equation}
\end{theorem}

\begin{proof}
The matrix identity follows by induction from
\(F_{n+1}=F_n+F_{n-1}\). The Perron root of \(F_{\rm av}\) is \(\varphi\),
hence \(h_{\rm top}=\log\varphi\). Since the matrix is symmetric, the Parry
weights \cite{Parry1964} are proportional to the squared components of a Perron vector
\((1,\varphi)\), that is, to \((1,\varphi^2)\).
\end{proof}

The expressions \(\log3\) and \(\log\varphi\) are thus distinguished by their
genealogy. The first quantifies a uniform local choice among three
states; the second measures the growth of histories allowed by
sheet incidence. The areal--volumetric relation is a property of the
complete return and therefore appears with exponent two.
